\documentclass{article}

\usepackage{amsmath}
\usepackage{amssymb}
\usepackage{booktabs}

\begin{document}

\appendix
\section{Analytical Structure of Nonlinear IPO Extensions}

\subsection{Why this appendix matters}

This appendix clarifies the analytical structure behind the nonlinear and
regularized IPO extensions used in the empirical study. The main distinction is
not whether a predictor is nonlinear in the input features, but whether it is
linear in its trainable parameters. Polynomial and finite-kernel predictors are
nonlinear in the input features but linear in their trainable coefficients, so
their IPO objectives retain a quadratic structure. Neural networks, by
contrast, are nonlinear in their parameters and therefore introduce nonconvex
optimization. This distinction helps interpret the empirical results: poor
performance of a nonlinear model may come from optimization instability and
high-variance predictions rather than from nonlinearity itself.

\subsection{Closed-form MVO decision layer}

We consider the integrated prediction-and-optimization problem
\[
\min_{\theta}
\frac{1}{T}\sum_{t=1}^T
c\bigl(z_t^*(\hat y_t), y_t\bigr),
\qquad
\hat y_t = f(X_t;\theta),
\]
where the downstream mean-variance portfolio decision is
\[
z_t^*(\hat y_t)
=
\arg\min_{z\in\mathcal S}
\left\{
-z^\top \hat y_t
+
\frac{\delta}{2}z^\top V_t z
\right\}.
\]

In the equality-constrained case
\[
\mathcal S = \{z: \mathbf 1^\top z = 1\},
\]
the optimal portfolio has the closed form
\[
z_t^*(\hat y_t)
=
\frac{1}{\delta}
V_t^{-1}
\left(
\hat y_t - \lambda_t \mathbf 1
\right),
\]
where
\[
\lambda_t
=
\frac{
\mathbf 1^\top V_t^{-1}\hat y_t - \delta
}{
\mathbf 1^\top V_t^{-1}\mathbf 1
}.
\]

Equivalently, the MVO solution is affine in the predicted return vector:
\[
z_t^*(\hat y_t)
=
B_t \hat y_t + a_t,
\]
where
\[
B_t
=
\frac{1}{\delta}
\left[
V_t^{-1}
-
\frac{
V_t^{-1}\mathbf 1 \mathbf 1^\top V_t^{-1}
}{
\mathbf 1^\top V_t^{-1}\mathbf 1
}
\right],
\]
and
\[
a_t
=
\frac{
V_t^{-1}\mathbf 1
}{
\mathbf 1^\top V_t^{-1}\mathbf 1
}.
\]

The key implication is that the MVO optimizer \(z_t^*(\hat y_t)\) is
closed-form for all predictors because it depends only on the predicted return
vector \(\hat y_t\). However, the predictor parameter estimator is
closed-form only when \(f(X_t;\theta)\) is linear in the trainable parameters
and the penalty is smooth.

\subsection{Parameter-linear and regularized IPO models}

Whenever the predictor is linear in its trainable parameters, the IPO objective
remains quadratic in those parameters after substituting the affine MVO
decision rule. This section derives the resulting structure for Ridge, Lasso,
Elastic Net, polynomial features, and finite-kernel predictors.

\subsubsection{Ridge IPO}

Consider a linear predictor
\[
\hat y_t = X_t \beta,
\]
where \(X_t \in \mathbb R^{n_t \times d}\), \(\beta \in \mathbb R^d\), and
\(\hat y_t \in \mathbb R^{n_t}\). Since the MVO solution is affine in
\(\hat y_t\), we have
\[
z_t^*(\hat y_t)
=
B_t X_t\beta + a_t.
\]

The realized MVO cost is
\[
c(z_t^*, y_t)
=
-(z_t^*)^\top y_t
+
\frac{\delta}{2}(z_t^*)^\top V_t z_t^*.
\]

Substituting \(z_t^* = B_tX_t\beta + a_t\), we obtain
\[
c_t(\beta)
=
-\left(B_tX_t\beta + a_t\right)^\top y_t
+
\frac{\delta}{2}
\left(B_tX_t\beta + a_t\right)^\top
V_t
\left(B_tX_t\beta + a_t\right).
\]

Ignoring terms independent of \(\beta\), the objective can be written as
\[
\mathcal L(\beta)
=
\frac{1}{2}\beta^\top H\beta
-
g^\top \beta
+
\frac{\lambda}{2}\|\beta\|_2^2,
\]
where
\[
H
=
\delta
\sum_{t=1}^T
X_t^\top B_t^\top V_t B_t X_t,
\]
and
\[
g
=
\sum_{t=1}^T
X_t^\top B_t^\top
\left(
y_t - \delta V_t a_t
\right).
\]

The ridge-regularized IPO estimator satisfies the first-order condition
\[
(H+\lambda I)\beta = g.
\]

Therefore, the closed-form ridge IPO solution is
\[
\boxed{
\hat\beta_{\mathrm{ridge\ IPO}}
=
(H+\lambda I)^{-1}g.
}
\]

\subsubsection{Lasso IPO}

The Lasso IPO objective adds an \(L_1\) penalty to the IPO loss:
\[
\mathcal L(\beta)
=
\frac{1}{2}\beta^\top H\beta
-
g^\top \beta
+
\lambda \|\beta\|_1.
\]

Equivalently,
\[
\hat\beta_{\mathrm{lasso\ IPO}}
=
\arg\min_{\beta}
\left\{
\frac{1}{2}\beta^\top H\beta
-
g^\top \beta
+
\lambda \|\beta\|_1
\right\}.
\]

Because the \(L_1\) norm is not differentiable at zero, the estimator generally
does not admit a simple closed-form expression like ridge regression. The
optimality condition is the subgradient condition
\[
0 \in H\hat\beta - g + \lambda \partial \|\hat\beta\|_1.
\]

Componentwise, this means
\[
(H\hat\beta - g)_j
+
\lambda s_j
=
0,
\]
where
\[
s_j =
\begin{cases}
\operatorname{sign}(\hat\beta_j), & \hat\beta_j \neq 0,\\
[-1,1], & \hat\beta_j = 0.
\end{cases}
\]

Thus, the Lasso IPO estimator is characterized by the KKT conditions:
\[
\boxed{
\begin{cases}
(H\hat\beta - g)_j = -\lambda \operatorname{sign}(\hat\beta_j),
& \hat\beta_j \neq 0,\\[4pt]
|(H\hat\beta - g)_j| \leq \lambda,
& \hat\beta_j = 0.
\end{cases}
}
\]

Only in the special case where \(H\) is diagonal does the solution reduce to
coordinatewise soft-thresholding:
\[
\hat\beta_j
=
\frac{1}{H_{jj}}
S(g_j,\lambda),
\]
where
\[
S(g_j,\lambda)
=
\operatorname{sign}(g_j)
\max\{|g_j|-\lambda,0\}.
\]

In the general IPO setting, \(H\) is not diagonal, so Lasso IPO is solved by
coordinate descent or proximal gradient methods rather than by a closed-form
matrix inverse.

\subsubsection{Elastic Net IPO}

The Elastic Net IPO objective is
\[
\mathcal L(\beta)
=
\frac{1}{2}\beta^\top H\beta
-
g^\top \beta
+
\lambda_1\|\beta\|_1
+
\frac{\lambda_2}{2}\|\beta\|_2^2.
\]

Combining the quadratic terms gives
\[
\mathcal L(\beta)
=
\frac{1}{2}\beta^\top (H+\lambda_2 I)\beta
-
g^\top \beta
+
\lambda_1\|\beta\|_1.
\]

Therefore,
\[
\hat\beta_{\mathrm{EN\ IPO}}
=
\arg\min_{\beta}
\left\{
\frac{1}{2}\beta^\top (H+\lambda_2 I)\beta
-
g^\top \beta
+
\lambda_1\|\beta\|_1
\right\}.
\]

The KKT condition is
\[
0
\in
(H+\lambda_2 I)\hat\beta
-
g
+
\lambda_1 \partial\|\hat\beta\|_1.
\]

Componentwise,
\[
\boxed{
\begin{cases}
\bigl((H+\lambda_2 I)\hat\beta - g\bigr)_j
=
-\lambda_1\operatorname{sign}(\hat\beta_j),
& \hat\beta_j\neq 0,\\[4pt]
\left|
\bigl((H+\lambda_2 I)\hat\beta - g\bigr)_j
\right|
\leq \lambda_1,
& \hat\beta_j=0.
\end{cases}
}
\]

As with Lasso IPO, the Elastic Net IPO estimator generally does not have a
closed-form matrix-inverse solution because of the non-differentiability of the
\(L_1\) term. It is typically solved by coordinate descent, proximal gradient,
or subgradient-based methods. When \(\lambda_1=0\), Elastic Net reduces to
Ridge IPO and admits the closed form
\[
\hat\beta
=
(H+\lambda_2 I)^{-1}g.
\]

\subsubsection{Polynomial IPO}

A polynomial predictor uses a nonlinear feature map
\[
\phi(x)
=
\left[
x,\;
x^{\circ 2},\;
\{x_jx_k\}_{j<k}
\right],
\]
where \(x^{\circ 2}\) denotes elementwise squares and
\(\{x_jx_k\}_{j<k}\) denotes pairwise interaction terms.

For a cross-section at time \(t\), let
\[
\Phi_t = \phi(X_t)
\]
denote the transformed feature matrix. The polynomial predictor is
\[
\hat y_t = \Phi_t \beta.
\]

Although the predictor is nonlinear in the original input features \(X_t\), it
is linear in the parameter vector \(\beta\). Therefore, using
\[
z_t^*(\hat y_t) = B_t\hat y_t + a_t,
\]
we obtain
\[
z_t^*
=
B_t\Phi_t\beta + a_t.
\]

The realized MVO cost becomes
\[
c_t(\beta)
=
-\left(B_t\Phi_t\beta+a_t\right)^\top y_t
+
\frac{\delta}{2}
\left(B_t\Phi_t\beta+a_t\right)^\top
V_t
\left(B_t\Phi_t\beta+a_t\right).
\]

Collecting the quadratic and linear terms across \(t\), define
\[
H_{\phi}
=
\delta
\sum_{t=1}^T
\Phi_t^\top B_t^\top V_t B_t \Phi_t,
\]
and
\[
g_{\phi}
=
\sum_{t=1}^T
\Phi_t^\top B_t^\top
\left(
y_t - \delta V_t a_t
\right).
\]

Then the polynomial IPO objective can be written as
\[
\mathcal L(\beta)
=
\frac{1}{2}\beta^\top H_{\phi}\beta
-
g_{\phi}^\top \beta.
\]

The first-order condition is
\[
H_{\phi}\beta = g_{\phi}.
\]

Therefore, the closed-form polynomial IPO solution is
\[
\boxed{
\hat\beta_{\mathrm{poly\ IPO}}
=
H_{\phi}^{-1}g_{\phi}.
}
\]

If an \(L_2\) penalty is added, the ridge-polynomial IPO solution becomes
\[
\boxed{
\hat\beta_{\mathrm{poly\ ridge\ IPO}}
=
(H_{\phi}+\lambda I)^{-1}g_{\phi}.
}
\]

\subsubsection{Finite-kernel IPO}

Consider a finite-dimensional kernel predictor with fixed anchor points
\(\{u_m\}_{m=1}^M\). Define the kernel feature matrix
\[
K_t
=
\left[
K(x_{i,t}, u_m)
\right]_{i=1,\ldots,n_t;\;m=1,\ldots,M}.
\]

The kernel predictor is
\[
\hat y_t = K_t \alpha + b\mathbf 1.
\]

After absorbing the intercept into an augmented kernel matrix
\[
\widetilde K_t =
\begin{bmatrix}
K_t & \mathbf 1
\end{bmatrix},
\qquad
\widetilde \alpha =
\begin{bmatrix}
\alpha\\
b
\end{bmatrix},
\]
we can write
\[
\hat y_t = \widetilde K_t\widetilde \alpha.
\]

Since the predictor is linear in the parameter vector
\(\widetilde\alpha\), the MVO solution becomes
\[
z_t^*
=
B_t\widetilde K_t\widetilde\alpha + a_t.
\]

The realized IPO objective is therefore quadratic in
\(\widetilde\alpha\):
\[
\mathcal L(\widetilde\alpha)
=
\frac{1}{2}
\widetilde\alpha^\top H_K \widetilde\alpha
-
g_K^\top \widetilde\alpha,
\]
where
\[
H_K
=
\delta
\sum_{t=1}^T
\widetilde K_t^\top B_t^\top V_t B_t \widetilde K_t,
\]
and
\[
g_K
=
\sum_{t=1}^T
\widetilde K_t^\top B_t^\top
\left(
y_t - \delta V_t a_t
\right).
\]

With kernel ridge regularization on the non-intercept parameters, let
\[
D =
\begin{bmatrix}
I_M & 0\\
0 & 0
\end{bmatrix}.
\]

The regularized kernel IPO objective is
\[
\mathcal L_{\mathrm{KRR}}(\widetilde\alpha)
=
\frac{1}{2}
\widetilde\alpha^\top H_K \widetilde\alpha
-
g_K^\top \widetilde\alpha
+
\frac{\lambda}{2}
\widetilde\alpha^\top D\widetilde\alpha.
\]

The first-order condition is
\[
(H_K+\lambda D)\widetilde\alpha = g_K.
\]

Thus, the finite-dimensional kernel ridge IPO estimator has the closed form
\[
\boxed{
\hat{\widetilde\alpha}_{\mathrm{kernel\ IPO}}
=
(H_K+\lambda D)^{-1}g_K.
}
\]

This closed form relies on the anchor points being fixed. If the kernel
parameters, such as the RBF bandwidth \(\gamma\), are also trained, then the
problem is no longer linear in the parameters and generally requires
gradient-based optimization.

\subsection{Fully nonlinear IPO models}

\subsubsection{Neural Network IPO}

A single-hidden-layer neural network predictor takes the form
\[
\hat y_t
=
f(X_t;\theta)
=
W_2 \sigma(W_1X_t + b_1) + b_2,
\]
where
\[
\theta = \{W_1,b_1,W_2,b_2\}.
\]

The MVO solution is still affine in the predicted return vector:
\[
z_t^*
=
B_t f(X_t;\theta) + a_t.
\]

Therefore, the IPO objective is
\[
\mathcal L(\theta)
=
\frac{1}{T}
\sum_{t=1}^T
\left[
-
\left(B_t f(X_t;\theta)+a_t\right)^\top y_t
+
\frac{\delta}{2}
\left(B_t f(X_t;\theta)+a_t\right)^\top
V_t
\left(B_t f(X_t;\theta)+a_t\right)
\right].
\]

Unlike the polynomial or finite-kernel cases, \(f(X_t;\theta)\) is nonlinear in
\(\theta\). As a result, \(\mathcal L(\theta)\) is generally nonconvex and does
not admit a closed-form estimator.

The neural-network IPO estimator is therefore defined as
\[
\boxed{
\hat\theta_{\mathrm{NN\ IPO}}
=
\arg\min_{\theta}
\mathcal L(\theta)
}
\]
and is solved by gradient-based optimization.

The gradient is computed through the differentiable MVO layer:
\[
\nabla_{\theta}\mathcal L(\theta)
=
\sum_{t=1}^T
\left(
\frac{\partial c_t}{\partial z_t^*}
\frac{\partial z_t^*}{\partial \hat y_t}
\frac{\partial f(X_t;\theta)}{\partial \theta}
\right).
\]

Since
\[
z_t^*(\hat y_t) = B_t\hat y_t + a_t,
\]
we have
\[
\frac{\partial z_t^*}{\partial \hat y_t} = B_t.
\]

Therefore,
\[
\nabla_{\theta}\mathcal L(\theta)
=
\sum_{t=1}^T
\left[
\left(
-y_t+\delta V_t z_t^*
\right)^\top
B_t
\frac{\partial f(X_t;\theta)}{\partial \theta}
\right],
\]
which can be evaluated by automatic differentiation. This is the approach used
for neural-network IPO.

\subsection{Summary of analytical structure}

Table \ref{tab:closed-form-summary} summarizes which model classes preserve a
closed-form or convex IPO structure.

\begin{table}[h]
\centering
\begin{tabular}{lccc}
\toprule
Model & Nonlinear in inputs? & Linear in parameters? & Closed-form IPO estimator? \\
\midrule
Linear & No & Yes & Yes \\
Ridge & No & Yes & Yes \\
Lasso & No & Yes & No, due to \(L_1\) nonsmoothness \\
Elastic Net & No & Yes & No, unless \(\lambda_1=0\) \\
Polynomial & Yes & Yes & Yes \\
Finite Kernel Ridge & Yes & Yes, if anchors fixed & Yes \\
Neural Network & Yes & No & No \\
\bottomrule
\end{tabular}
\caption{Closed-form availability for IPO extensions. The key distinction is
not whether the predictor is nonlinear in the inputs, but whether it is linear
in the trainable parameters.}
\label{tab:closed-form-summary}
\end{table}

The key point is that the MVO optimizer \(z^*(\hat y_t)\) is closed-form for
all predictors because it depends only on the predicted return vector
\(\hat y_t\). However, the predictor parameter estimator is closed-form only
when \(f(X_t;\theta)\) is linear in the trainable parameters and the penalty is
smooth. This distinction explains why ridge, polynomial, and finite-kernel
models have closed-form IPO structure, while Lasso and Elastic Net require
nonsmooth convex optimization and neural networks require gradient-based
nonconvex optimization.

\subsection{Connection to implementation and empirical results}

The derivations above clarify the analytical structure of the IPO objective.
In the implementation, however, we use a unified training interface. The OLS
plug-in baselines use closed-form least squares when available and sklearn
coordinate descent for Lasso and Elastic Net. The IPO variants are trained
end-to-end by backpropagating realized MVO cost through the closed-form
equality-constrained MVO layer. Thus, the closed form is used operationally for
\(z^*(\hat y_t)\) in every IPO model, while the predictor parameters are updated
by gradient descent except in the OLS plug-in baselines.

The theoretical classification also helps explain the empirical ranking.
Polynomial and finite-kernel predictors are flexible but high-dimensional; even
though their IPO objectives retain a tractable structure, the lifted feature
space can amplify estimation noise in \(\hat y_t\), which is then magnified by
the MVO layer into extreme portfolio weights. In contrast, Lasso and Elastic
Net impose sparsity directly on the predictor, reducing the variance of
\(\hat y_t\) and producing more stable portfolio decisions.

Table \ref{tab:implementation-mapping} maps the theoretical structure to the
implemented model classes.

\begin{table}[h]
\centering
\begin{tabular}{lccc}
\toprule
Model & Predictor class & Analytical structure & Implementation \\
\midrule
Linear/Ridge
& Parameter-linear
& Quadratic IPO objective
& OLS closed-form; IPO Adam \\

Lasso/Elastic Net
& Parameter-linear + \(L_1\)
& Convex but nonsmooth
& sklearn / Adam with \(L_1\) penalty \\

Polynomial
& Nonlinear features, parameter-linear
& Quadratic after feature expansion
& Feature lift + linear layer \\

Finite Kernel
& Fixed-anchor nonlinear features
& Quadratic in kernel weights
& Fixed anchors + trainable weights \\

Neural Network
& Nonlinear in parameters
& Nonconvex IPO objective
& Backprop through MVO layer \\
\bottomrule
\end{tabular}
\caption{How the theoretical IPO structure maps to the implemented model classes.}
\label{tab:implementation-mapping}
\end{table}

\subsection{Covariance estimation}

At each monthly rebalance date \(t\), the covariance matrix \(V_t\) is estimated
from the trailing 60 daily returns of the assets in the current cross-section.
Because the cross-section is large relative to the estimation window, we use a
Ledoit--Wolf linear shrinkage estimator toward a scaled identity matrix. The
resulting daily covariance estimate is multiplied by 21 so that it is expressed
in the same monthly units as the realized next-month return \(y_t\).

\subsection{Limitations of the analytical structure}

The derivations characterize the optimization structure conditional on a fixed
covariance estimate \(V_t\). They do not solve the joint problem of learning
both expected returns and the covariance matrix, nor do they incorporate
long-only, leverage, or turnover constraints. These extensions would generally
destroy the simple affine MVO layer and require either differentiable quadratic
programming layers or approximate projection methods.

\end{document}